% ===================================================================
%  TAULA N.º 7 — El Poble a l'hivern. — La Masia a la primavera.
%  Traduction 2026 d'apres texto/io, controlee sur texto/fr et
%  texto/es. Consigne : voir texto/ca/10-taula-01.tex.
% ===================================================================

%%K t07-noto-1 noto
\VUnotes{143.2mm}{%
(*) Per al detall d'un àpat, vegeu les taules 6 i 13.%
}

%%K t07-apar-1 apar
\VUsaut{4.0mm}
\VUcentre{13.2pt}{90}{SEGONA SÈRIE}
\VUsaut{3.0mm}
\VUfilet{16mm}
\VUsaut{5.6mm}
\VUtitre{15.0pt}{60}{TAULA N.º 7}
\VUsaut{3.0mm}

%%K t07-tit-1 sub
\VUcentre{12.0pt}{0}{\textit{Primera escena.}}
\VUsaut{3.0mm}

%%K t07-tit-2 sub
\VUpk{10.2pt}{40}{El poble a l'hivern.}
\VUsaut{4.0mm}

%%K t07-c1-01-1 p
1. --- Durant les vacances de Cap d'Any vaig acompanyar el meu pare a Vilanova,
un poblet on tenia alguns afers per tractar.

%%K t07-c1-02-1 p
2. --- Vam agafar la \VUgras{diligència} \textsuperscript{(11)} que fa el servei
entre la ciutat i aquell poble; però com que feia molt de fred, aquell viatge en
un vehicle tan poc còmode no va ser gens agradable.

%%K t07-c1-03-1 p
3. --- Per això vam sentir veritable plaer quan vam veure que el \VUgras{cotxer}
aturava el seu carruatge a la plaça, davant de l'\VUgras{hostal}
\textsuperscript{(2)} de l'\textit{Ós Blanc}. L'\VUgras{hostaler}
\textsuperscript{(4)} ens esperava al llindar de casa seva, fumant tranquil·lament
la seva \VUgras{pipa}.

%%K t07-c1-03-2 p
Ens va convidar a entrar i no em vaig fer pregar per instal·lar-me davant del foc
de sarments que espurnejava a la llar. Aleshores em reia de l'aspre \VUgras{vent
del nord} que feia grinyolar el vell \VUgras{rètol} \textsuperscript{(3)} i
s'enduia en negres remolins el \VUgras{fum} \textsuperscript{(1)} que sortia de la
xemeneia.

%%K t07-c1-04-1 p
4. --- Era l'hora de dinar. Sense esperar més, vam fer honor al dinar senzill
però saborós que ens van servir (*).

%%K t07-tit-3 sub
\VUpk{10.2pt}{40}{Unes quantes visites.}
\VUsaut{3.0mm}

%%K t07-c1-05-1 p
5. --- Després de dinar vaig acompanyar el meu pare, que és agent d'assegurances,
a casa dels seus clients. Primer vam visitar el \VUgras{ferrer}
\textsuperscript{(5)}, que viu vora l'hostal. Estava ferrant un cavall. Vaig
admirar la força del seu ajudant, que subjectava fermament contra el seu davantal
de cuir el \VUgras{casc} del cavall mentre el ferrer fixava la \VUgras{ferradura}
\textsuperscript{(6)} a forts cops de martell.

%%K t07-c1-06-1 p
6. --- Després vam anar a casa del \VUgras{sabater} \textsuperscript{(7)} del
poble, l'oncle Jaume. Encara que tenia per rètol una bota \VUgras{magnífica}, es
conformava a posar soles a un parell de sabates velles i foradades.

%%K t07-c1-06-2 p
El bon home ens va dir rient: «A la ciutat era “sabater” i no tenia clients. Aquí
sóc “sabater de vell” i hi ha feina de sobres».

%%K t07-c1-07-1 p
7. --- Ens va parlar del seu veí el \VUgras{barber} \textsuperscript{(8)}, la
clientela del qual està descontenta. Les seves \VUgras{navalles} són sempre
esmussades i les seves \VUgras{tisores} no tallen mai. Però com que és l'únic
barber del poble, no hi ha més remei que deixar-se afaitar i tallar el cabell per
ell. A més és un home avar i antipàtic.

%%K t07-c1-08-1 p
8. --- Per sort, els altres \VUgras{veïns} no s'hi assemblen. El \VUgras{carreter}
\textsuperscript{(9)}, per exemple, és un home excel·lent, treballador i
servicial. Fa gust de portar-li un cotxe a arreglar. La seva feina és sempre ben
acabada.

%%K t07-c1-09-1 p
9. --- L'oncle Jaume tampoc no es queixava del \VUgras{baster}
\textsuperscript{(10)}, a qui de vegades ajuda a cosir els \VUgras{arreus} quan la
feina apressa.

%%K t07-c1-09-2 p
El meu pare li va preguntar per la seva família. «Tots estan bé. La meva néta és
a l'\VUgras{escola} \textsuperscript{(12)}. La seva \VUgras{mestra}
\textsuperscript{(13)} n'està molt contenta; és bona \VUgras{alumna}
\textsuperscript{(14)}. No és com aquell múrria del meu fill; només pensa a
jugar. El \VUgras{mestre} \textsuperscript{(15)} no aconsegueix ensenyar-li res».

%%K t07-tit-4 sub
\VUsaut{3.0mm}
\VUpk{10.2pt}{40}{Xafarderies del poble.}
\VUsaut{3.0mm}

%%K t07-c1-10-1 p
10. --- El bon home ens va explicar després les novetats del poble. Anaven a
construir una escola nova, perquè la que hi havia instal·lada a la \VUgras{casa de
la vila} s'havia quedat molt petita. Altrament era fàcil, perquè n'hi havia prou
de comprar una part del jardí del \VUgras{moliner}. Això li aniria molt bé al
pobre home. D'ençà que van arribar els freds, les \VUgras{moles} del \VUgras{molí}
\textsuperscript{(16)} ja no podien moldre el blat, perquè la \VUgras{roda}
\textsuperscript{(17)} estava presa en el \VUgras{gel} \textsuperscript{(25)} del
\VUgras{rierol} \textsuperscript{(23)}.

%%K t07-c1-11-1 p
11. --- «I parlant de construccions, han vist la nostra nova \VUgras{església}
\textsuperscript{(18)}? Està molt pintoresca sota el seu mantell de neu. Els
\VUgras{corbs} \textsuperscript{(19)}, que ja no reconeixen el seu vell campanar,
es posen tristament al gran arbre del \VUgras{cementiri}
\textsuperscript{(20)}».

%%K t07-c1-12-1 p
12. --- La idea del cementiri va recordar al sabater les persones que el meu pare
havia conegut i que havien mort l'any passat. «Quantes \VUgras{tombes}
\textsuperscript{(21)}, bon senyor meu, s'han afegit a les altres sota el
\VUgras{xiprer} \textsuperscript{(22)}! La mort es va endur el nostre vell notari.
Tot el poble va assistir al seu \VUgras{enterrament}».

%%K t07-c1-13-1 p
13. --- «Vet aquí el seu successor, patinant al rierol. Va venir fa una estona
perquè li arreglés la corretja del seu \VUgras{patí} \textsuperscript{(26)}, que
s'havia trencat».

%%K t07-c1-14-1 p
14. --- «Vet aquí també, a la vora del rierol, el nou \VUgras{guarda rural}
\textsuperscript{(27)}. És un antic gendarme i el terror dels trinxeraires del
poble. Fugen tan bon punt veuen les seves \VUgras{polaines}
\textsuperscript{(29)} i el seu \VUgras{quepis} \textsuperscript{(28)}».

%%K t07-c1-15-1 p
15. --- «Ah! Vet aquí l'oncle Josep, que porta un \VUgras{carro de feixos de
llenya} \textsuperscript{(30)} per al forner. --- És veritat, va dir el meu pare,
reconec el seu tir de \VUgras{mules} \textsuperscript{(31)}. He de parlar amb ell
i aprofitaré l'ocasió. A reveure, oncle Jaume».

%%K t07-c1-16-1 p
16. --- Just quan sortíem, els nens del poble deixaven l'escola i es llançaven
corrents al \VUgras{lliscador} \textsuperscript{(33)}, amb risc de caure i
trencar-se un membre en la \VUgras{caiguda} \textsuperscript{(34)}. Un d'ells va
agafar el seu \VUgras{trineu} \textsuperscript{(32)} a casa del carreter, hi va
asseure un dels seus companys i va partir corrents.

%%K t07-c1-17-1 p
17. --- D'altres es van precipitar cap a un \VUgras{munt de neu}
\textsuperscript{(37)} vora el \VUgras{pont} \textsuperscript{(40)} i van començar
a bombardejar el \VUgras{ninot de neu} \textsuperscript{(39)} que havien fet el dia
abans i a qui havien posat un barret, una pipa i una cartera d'escola en
bandolera.

%%K t07-c1-18-1 p
18. --- Poc els importaven el vent gelat i el fred. La majoria ni tan sols tenia
\VUgras{bufanda} \textsuperscript{(35)} i, per entrar en calor després de la
batalla de \VUgras{boles de neu} \textsuperscript{(38)}, els bastava un grapat de
\VUgras{castanyes} \textsuperscript{(42)} molt calentes comprades a la tia Jaumeta,
la \VUgras{venedora}, que tremola de fred sota el seu \VUgras{xal}
\textsuperscript{(43)} massa prim.

%%K t07-tit-5 sub
\VUsaut{3.0mm}
\VUcentre{12.0pt}{0}{\textit{Segona escena.}}
\VUsaut{3.0mm}

%%K t07-tit-6 sub
\VUpk{10.2pt}{40}{La casa de pagès a la primavera.}
\VUsaut{4.0mm}

%%K t07-c2-01-1 p
1. --- Malgrat tots els plaers de l'hivern, amb pregona joia saludem el retorn de
la \VUgras{primavera}.

%%K t07-c2-01-2 p
Quin quadre tan alegre és el corral d'una masia, al capvespre, al mes de maig!

%%K t07-c2-02-1 p
2. --- A l'horitzó, el \VUgras{sol} \textsuperscript{(1)} es pon a poc a poc,
inundant el \VUgras{camp} \textsuperscript{(3)} amb els seus \VUgras{raigs}
\textsuperscript{(2)} de porpra. El diligent \VUgras{llaurador}
\textsuperscript{(6)} s'afanya a llaurar el seu \VUgras{camp}
\textsuperscript{(4)}.

%%K t07-c2-02-2 p
Amb la seva \VUgras{arada} \textsuperscript{(5)}, tirada per un vigorós
\VUgras{cavall} \textsuperscript{(7)}, traça el \VUgras{solc}
\textsuperscript{(8)} en què el \VUgras{sembrador} \textsuperscript{(9)} tirarà a
grapats la \VUgras{llavor} que donarà les espigues daurades.

%%K t07-c2-03-1 p
3. --- Els \VUgras{segadors} \textsuperscript{(11)}, amb les seves dalles, han
tallat l'herba del prat, els assecadors han assecat el \VUgras{fenc}
\textsuperscript{(10)} girant-lo amb les \VUgras{forques} \textsuperscript{(12)} i
els rasclets, i vet-los aquí de tornada.

%%K t07-c2-03-2 p
Uns van davant del pesat carro tirat per bous; porten la seva eina a l'espatlla.
D'altres, més peresosos, s'han instal·lat còmodament damunt el fenc olorós.

%%K t07-c2-04-1 p
4. --- En passar saluden amistosament el \VUgras{jardiner} \textsuperscript{(16)},
ocupat a l'\VUgras{hort de fruiters} \textsuperscript{(13)} a podar els
\VUgras{arbres fruiters} i empeltar les \VUgras{pereres} \textsuperscript{(14)}.
Les \VUgras{prunieres} \textsuperscript{(15)} són en flor i, a l'\VUgras{hort}
\textsuperscript{(17)} ben adobat, les verdures, cols i enciams ja són bells.

%%K t07-c2-04-2 p
Damunt la paret, un bell \VUgras{paó} \textsuperscript{(18)} desplega la cua
perquè admirin les seves magnífiques plomes.

%%K t07-tit-7 sub
\VUpk{10.2pt}{40}{El corral de la masia.}
\VUsaut{3.0mm}

%%K t07-c2-05-1 p
5. --- Les portes de la \VUgras{quadra} \textsuperscript{(19)} són obertes i es
veuen la menjadora i el \VUgras{jaç de palla} \textsuperscript{(23)} de l'\VUgras{euga}
\textsuperscript{(21)} que el \VUgras{mosso de quadra} \textsuperscript{(20)} acaba
de desenganxar. El \VUgras{poltre} \textsuperscript{(22)}, tan graciós amb les
seves potes llargues, segueix la seva mare fent bots.

%%K t07-c2-06-1 p
6. --- Vora el \VUgras{pou} \textsuperscript{(29)}, les \VUgras{vaques}
\textsuperscript{(25)} van a beure a l'\VUgras{abeurador} \textsuperscript{(26)} de
fusta que els serveix de pica. El \VUgras{vedell} \textsuperscript{(24)} ho
aprofita per mamar, i el \VUgras{bou} \textsuperscript{(27)}, olorant la bona olor
del farratge, mugeix alegrement i alça amb orgull el cap de poderoses
\VUgras{banyes} \textsuperscript{(28)}.

%%K t07-c2-07-1 p
7. --- Tots treballen. La \VUgras{criada} \textsuperscript{(32)} sosté amb la mà la
\VUgras{corda} \textsuperscript{(31)} del pou. Treu aigua amb un \VUgras{cubell}
\textsuperscript{(30)} per donar de beure al bestiar. Vora el \VUgras{graner}
\textsuperscript{(33)}, un home rastella els brins de palla, i davant de la porta
de la \VUgras{casa} \textsuperscript{(34)}, una vella \VUgras{filadora}
\textsuperscript{(35)}, amb la seva filosa i el seu \VUgras{fus}, fila el
\VUgras{lli} o el \VUgras{cànem} com en altres temps.

%%K t07-tit-8 sub
\VUsaut{3.0mm}
\VUpk{10.2pt}{40}{El pagès.}
\VUsaut{3.0mm}

%%K t07-c2-08-1 p
8. --- El \VUgras{pagès} \textsuperscript{(36)} dirigeix tota aquesta feina i dóna
ordres als seus mossos. Mana recollir el \VUgras{rascle} \textsuperscript{(40)} i
el \VUgras{pic} \textsuperscript{(39)} que s'han oblidat vora la \VUgras{caseta}
\textsuperscript{(38)} del \VUgras{gos} \textsuperscript{(37)}.

%%K t07-c2-09-1 p
9. --- Els seus crits espanten un \VUgras{colom} \textsuperscript{(41)}, que fuig a
tot vol cap al \VUgras{colomar} \textsuperscript{(42)}, i les \VUgras{gallines}
\textsuperscript{(43)}, que s'afanyen a tornar al \VUgras{galliner}
\textsuperscript{(44)}.

%%K t07-c2-10-1 p
10. --- Davant del \VUgras{corral d'ovelles} \textsuperscript{(48)}, heus aquí el
vell \VUgras{pastor} \textsuperscript{(45)} que torna amb el seu \VUgras{ramat}
\textsuperscript{(49)}. Amb el seu \VUgras{gaiato} \textsuperscript{(46)}, que no
el deixa mai, igual que la seva \VUgras{brusa} \textsuperscript{(47)} descolorida,
mena al corral \VUgras{marrans} \textsuperscript{(50)}, \VUgras{ovelles}
\textsuperscript{(53)}, \VUgras{anyells} \textsuperscript{(54)},
\VUgras{cabres} \textsuperscript{(51)} i \VUgras{cabrits}
\textsuperscript{(52)}. El seu fidel \VUgras{gos} \textsuperscript{(55)} Argos
tanca la marxa i anima amb els seus lladrucs els endarrerits.

%%K t07-c2-11-1 p
11. --- Tot aquest soroll i aquest moviment no torben la calma de la bona
\VUgras{vaca lletera} \textsuperscript{(56)}, que fa sonar la seva \VUgras{esquella}
\textsuperscript{(57)} mentre la munyen davant de la porta de l'\VUgras{estable}
\textsuperscript{(58)}.

%%K t07-c2-12-1 p
12. --- Sembla comprendre que ha de donar la seva llet perquè la \VUgras{pagesa}
\textsuperscript{(60)} pugui fer \VUgras{mantega} a la \VUgras{manteguera}
\textsuperscript{(61)} o els excel·lents \VUgras{formatges}
\textsuperscript{(64)} que escorren damunt la post posada damunt el
\VUgras{cossi} \textsuperscript{(63)}.

%%K t07-c2-13-1 p
13. --- Tota la \VUgras{lleteria} \textsuperscript{(59)} la manté molt neta el
\VUgras{mosso de la masia} \textsuperscript{(65)}, que acaba de rentar els
\VUgras{càntirs de la llet} \textsuperscript{(62)} al gran cubell que porta a la
mà.

%%K t07-tit-9 sub
\VUsaut{3.0mm}
\VUpk{10.2pt}{40}{El corral de les aus.}
\VUsaut{3.0mm}

%%K t07-c2-14-1 p
14. --- Amb els seus gruixuts esclops camina amb una mica de matusseria, i així
fa fugir el \VUgras{gall dindi} \textsuperscript{(68)}, les \VUgras{ànedes}
\textsuperscript{(66)}, els \VUgras{ànecs} \textsuperscript{(67)} i les
\VUgras{oques} \textsuperscript{(69)}, que obren molt el \VUgras{bec}
\textsuperscript{(70)} i llancen crits inquiets.

%%K t07-c2-14-2 p
El \VUgras{gall} \textsuperscript{(71)}, més bel·licós, alça la seva
\VUgras{cresta} \textsuperscript{(72)} vermella, sacseja les plomes de la seva
\VUgras{cua} \textsuperscript{(73)} i llança un sonor «quiquiriquic».

%%K t07-c2-15-1 p
15. --- Una \VUgras{lloca} \textsuperscript{(74)} picoteja al \VUgras{femer}
\textsuperscript{(81)} amb els seus \VUgras{pollets} \textsuperscript{(75)}. El
\VUgras{porcell} \textsuperscript{(78)} corre cap a la seva mare, l'enorme
\VUgras{truja} \textsuperscript{(77)} que veiem vora la cort.

%%K t07-c2-16-1 p
16. --- A les seves gàbies, els \VUgras{conills} \textsuperscript{(79)} mengen amb
avidesa l'\VUgras{herba} \textsuperscript{(80)} tendra que els porta la filla del
pagès. Joaneta estima molt els seus conills i, cada dia, després d'anar a buscar
els \VUgras{ous} \textsuperscript{(76)} frescos al galliner, ve a visitar els seus
amiguets.
\VUsaut{6.0mm}
\VUfilet{22mm}
